\section{Herramientas síncronas vs. asíncronas: las tres eras de la programación con IA}

\subsection{Conferencia invitada: Silas Alberti (Cognition/Devin)}

Silas Alberti, director de investigación de Cognition, compartió la trayectoria de evolución de las herramientas de programación con IA y las mejores estrategias para usarlas.

\subsection{Las tres eras de las herramientas de programación con IA}

\begin{enumerate}
    \item \textbf{Code Completion} ($\sim$10\% de mejora en eficiencia): GitHub Copilot acelera la codificación
    \item \textbf{IDE Automation} ($\sim$20\% de mejora en eficiencia): el AI IDE logra que una sola persona complete la tarea
    \item \textbf{AI Software Engineer} (6--12$\times$ de mejora en eficiencia): los Agentes en la nube logran flujos de trabajo paralelos
\end{enumerate}

\begin{importantbox}{De 10\% a 6--12$\times$: de cambio cuantitativo a cambio cualitativo}
De code completion a cloud agent, la mejora en eficiencia pasa de un nivel porcentual a un nivel de múltiplos. La diferencia clave es que el cloud agent logra el \textbf{paralelismo}: un desarrollador puede administrar varios Agentes al mismo tiempo para ejecutar tareas distintas.
\end{importantbox}

\subsection{Herramientas síncronas vs. asíncronas}

\textbf{Herramientas síncronas} (Sync):
\begin{itemize}
    \item Un solo hilo, con el ser humano en el bucle (human-in-the-loop)
    \item La atención se concentra en una sola tarea
    \item El Agent trabaja de 20 segundos a 1.5 minutos
    \item Representantes: Cursor, Windsurf
\end{itemize}

\textbf{Herramientas asíncronas} (Async):
\begin{itemize}
    \item Múltiples hilos; el ser humano delega la tarea y luego cambia su atención a otra cosa
    \item El Agent trabaja de 10 minutos a varias horas
    \item Representante: Devin
\end{itemize}

\begin{warningbox}{La trampa del «Semi-Async»}
Hay que tener especial cuidado en el rango «semiasíncrono» de 30 segundos a 3 minutos: es demasiado lento para mantener el flujo (flow) y, al mismo tiempo, demasiado corto para justificar el cambio a otra tarea. La estrategia correcta es hacer que la herramienta sea más rápida para mantener el flujo, o bien invertir más tiempo a cambio de un nivel de inteligencia más alto.
\end{warningbox}

\subsection{El mejor flujo de trabajo de programación en 2025}

La mejor práctica que propone Silas asigna las tres etapas a herramientas distintas:

\begin{itemize}
    \item \textbf{Planning} (planeación): herramientas síncronas — usar DeepWiki y Codemaps para entender la base de código
    \item \textbf{Coding} (codificación): herramientas asíncronas — delegar la ejecución a Devin
    \item \textbf{Testing} (pruebas): herramientas síncronas — probar y ajustar en Windsurf los resultados del Agent
\end{itemize}

\subsection{Perspectivas a futuro}

\begin{enumerate}
    \item Los ingenieros humanos se convierten en administradores de Agentes: las herramientas síncronas resuelven los problemas más difíciles y las asíncronas obtienen un apalancamiento de 10$\times$
    \item Las habilidades clave del futuro: \textbf{delegación y manejo de múltiples hilos}, \textbf{lectura de código}, \textbf{planeación, delimitación del alcance y diseño de arquitectura}
    \item Tendencia a largo plazo: las tres etapas de planeación, codificación y pruebas se volverán, de manera gradual, asíncronas y autónomas
\end{enumerate}

\subsection{Marco de evaluación para que un equipo adopte un AI IDE}

Si se traduce el contenido de la Semana 3 a criterios de decisión a nivel organizacional, las siguientes tres preguntas permiten evaluar rápidamente si vale la pena adoptar un AI IDE en toda la organización:

\begin{itemize}
    \item \textbf{¿Reduce de verdad el cambio de contexto?} Si los desarrolladores siguen teniendo que saltar entre una gran cantidad de documentos externos y plantillas de prompt, la herramienta no aporta una mejora sustancial.
    \item \textbf{¿Hace más transparente la base de código?} Un buen AI IDE amplifica la información estructurada que ya existe en el repo, en lugar de ocultar el desorden.
    \item \textbf{¿Cuenta con un mecanismo de reversión estable?} Cuando el Agent hace una modificación incorrecta, ¿puede el equipo apoyarse en pruebas, en el review y en un diff claro para recuperarse con rapidez?
\end{itemize}

\begin{knowledgebox}{Un AI IDE realmente viable en la práctica debe optimizar tanto a las personas como al repositorio}
Un AI IDE nunca es solo un problema de modelo. También exige una mejor organización del código, un sistema de evaluación más estable y una documentación del proyecto más clara. Precisamente por eso, el context engineering termina obligando a los equipos a reestructurar su forma de administrar el conocimiento.
\end{knowledgebox}

\subsection{Resumen de la sección}

Las herramientas de programación con IA evolucionaron de code completion a cloud agent y trajeron un salto en la eficiencia de 10\% a 6--12$\times$. Lo importante es dominar los escenarios de uso correctos de las herramientas síncronas y asíncronas, evitar la trampa del «semiasíncrono» y asignar la planeación, la codificación y las pruebas al tipo de herramienta más adecuado.
